\begin{description}
	\item[\texttt{model}ブロック] 尤度と事前分布を書くところ。確率分布からのサンプリングは変数を左に，$\sim$で\keytermL{確率分布}{かくりつぶんぷ}につなぎます。\texttt{normal}とか\texttt{cauchy}が確率分布の名前なのでRStudioではハイライト表示されているはず。一行の終わりはセミコロンで。複数のデータや変数，$Y_i$や$\sigma_i$は，iごとに変わるということをあらわしていますが，プログラム的には\verb|Y[i]|のように書いて，大括弧で変数を括って添字であることを表現します。データは7件あって，それぞれについて$i$さんの測定値$Y_i$と測定誤差$\sigma_i$があるので，\texttt{for}文で\verb|i|を繰り返しています。変数\verb|sig[i]|についての事前分布も同様に繰り返しています。ちなみに変数名\texttt{mu, sig, Y}などは任意で，予約語でなければ好きな名前にしてもらって構いません。
	\item[\texttt{parameters}ブロック] \texttt{model}ブロックで，変数\texttt{mu}や\texttt{sig}という任意に命名して使っていたわけですが，それは何なのかを宣言してやらねばなりません。これらは求めたいパラメータなのですから，このブロックに型と共に宣言するわけです。変数\verb|mu|は実数型，変数\verb|sig|も実数型ですが7つの要素を持つ配列，しかも負の数を取らないので下限はゼロですよ，という宣言をしているところになります。範囲の設定に\verb|<>|の記号を使っていること注意してください。配列については次の\texttt{data}ブロックで説明します。
	\item[\texttt{data}ブロック] 最後にデータブロックです。モデルの変数でもパラメータでもない，外部から与えられるものなので，変数名\verb|Y|で宣言します。宣言のときに1つの変数に複数の要素がある(\keytermL{配列}{はいれつ}といいます)ことを明示するため，\verb|array[size]|と書きます。この大括弧(\verb|[ ]| )の中身は整数で，今回は7件のデータなのでサイズは7としてあります。そしてその配列の数字が整数(Integer)なのか，実数(Real)なのか，といった型宣言をしたうえで，変数名を書きます\footnote{Stanの以前のバージョンでは\verb|Y[i]|や\verb|sig[i]|のように，大括弧(\verb|[ ]|)を後ろにつけて書いていました。しかしバージョン2.32.0はその表記は誤りとされ，ここで解説しているような\verb|array|表記がスタンダードになります。古い記法のままだと警告が出たり，RStudioも古いバージョンだと新しい記法に変更要請(赤い波線がコードについたり)が出たりしますが，今すぐ問題になるというものではありません。とはいえ，なるべく新しいバージョンを使ったほうがよいので，表記法も新しいものに心がけましょう。}。
\end{description}
